% Generated draft: requires mathematical and source review.
\documentclass[11pt]{article}
\usepackage[margin=25mm]{geometry}
\usepackage{fontspec}
\setmainfont{Latin Modern Roman}
\usepackage{amsmath,amssymb,mathtools,unicode-math}
\setmathfont{Latin Modern Math}
\usepackage{enumitem,needspace,xurl,hyperref}
\setlength{\parindent}{0pt}
\setlength{\parskip}{5pt}
\setlength{\emergencystretch}{3em}
\newcommand{\N}{\mathbb{N}}
\newcommand{\Z}{\mathbb{Z}}
\newcommand{\Q}{\mathbb{Q}}
\newcommand{\R}{\mathbb{R}}
\newcommand{\C}{\mathbb{C}}
\newcommand{\F}{\mathbb{F}}
\newcommand{\A}{\mathbb{A}}
\newcommand{\PP}{\mathbb{P}}
\newcommand{\E}{\mathbb{E}}
\newcommand{\eps}{\varepsilon}
\newcommand{\dd}{\,\mathrm{d}}
\begin{document}
\section*{UnsolvedMath ID: 1200034}
\textbf{Source problem code:} catalogue entry
\subsection{Definitions and mathematical statement}
Use the measure-theoretic Bernoulli-shift interpretation. Let $\Gamma$ be a countable discrete nonamenable group, where nonamenable means that no left-translation-invariant probability mean exists on the bounded real functions on $\Gamma$. For $j=2,3$, let $A_2=\{0,1\}$ and $A_3=\{0,1,2\}$, let $u_j$ give each symbol mass $1/j$, and put $X_j=A_j^\Gamma$ with product measure $\mu_j=u_j^\Gamma$. The left shift action is
\[
(\gamma\cdot x)(h)=x(\gamma^{-1}h)\qquad(\gamma,h\in\Gamma).
\]
A measurable factor from $X_2$ onto $X_3$ is a measurable map $\Phi:X_2\to X_3$, defined modulo null sets, such that
\[
\Phi_*\mu_2=\mu_3,\qquad
\Phi(\gamma\cdot x)=\gamma\cdot\Phi(x)
\quad\text{for every }\gamma\in\Gamma\text{ and }\mu_2\text{-almost every }x.
\]
Onto here is in the probability-space sense; no everywhere-defined set-theoretic surjection or continuity is required.

With these conventions, the question is whether such a factor exists for every countable nonamenable $\Gamma$. The original source says only ``nonamenable group'' and writes the two shift spaces without specifying their base measures or factor category. The countability restriction is explicit in the supplied resolution summary, and the displayed formulation states the usual uniform-product measurable interpretation. It does not silently assert a continuous-factor theorem, a biased-measure variant, or a result for uncountable groups.

The supplied summary labels the answer affirmative, naming the Ornstein--Weiss result for the free group $\mathbb F_2$ and results of Bowen, Seward, and Tucker-Drob for the general countable setting. This is a reported literature resolution, not a proof or independent verification supplied here.
\subsection{Short English statement}
For every countable nonamenable group, can the uniform two-symbol Bernoulli action measurably produce the uniform three-symbol Bernoulli action while preserving the group action and output distribution? The supplied summary says yes. The original does not specify measures, countability, or the factor category, so the usual measurable uniform-product interpretation is stated explicitly.
\subsection{Sources}
\begin{itemize}
\item Local source registry: \url{data/unsolvedmath/problems.json}, permanent UnsolvedMath ID $1200034$, source code catalogue entry, including its supplied solved classification and research summary. Dataset attribution: Ulam AI and UnsolvedMath Contributors; original-source attribution and rights are retained.
\item Abert, \emph{Some questions} (2010), Question 34: \url{https://www.renyi.hu/~abert/questions.pdf}. The supplied metadata reports manual transcription and July 29, 2026 access, with current-status and rights review flags.
\item The supplied summary attributes the free-group case to Ornstein--Weiss and the general countable case to Bowen and Seward, also naming Tucker-Drob. It supplies no paper titles, dates, theorem numbers, or separate links for these resolution claims; they were not independently checked here.
\end{itemize}
\end{document}
